% ===========================================================================
% Section 2: Historical Context
%
% Ported from the authoritative PDF at Plan/paper.pdf (copied to
% paper/paper_planned.pdf). Four subsections cover the railway system
% before 1960, the Larkin Plan, the implementation of closures and road
% expansion, and the scale-economies motivation for sector-specific
% analysis.
%
% Per Shapiro Robot-mode (tasks W1-W4), this section "states facts";
% interpretation of the findings belongs in later sections. The
% pre-existing text is already in Robot mode and has been preserved
% with only minor cleanups (line-break artifacts from PDF extraction).
%
% PLACEHOLDERS:
%   - [1 paragraph on freight evidence] at the end of 2.4 — pending
%     access to the Santa Fe Railroad freight data. Coauthor note:
%     "Frame as suggestive, not conclusive."
%
% Citation keys used:
%   baumgartnerpalazzo1969, fajgelbaumredding2022, keeling1993, larkin1962,
%   lopez_waddell2007, pahowka2005, potash1969, rock1987, whitaker1954
% All keys resolve in references.bib (verified entries, PR #89).
% ===========================================================================

\section{Historical Context}
\label{sec:history}

\subsection{Argentina's Railway System and Fiscal Crisis}
\label{ssec:railway_history}

Argentina's railway network was one of the largest in the world.
Construction began in 1857 and expanded rapidly through a combination
of national capital, foreign investment—predominantly British and
French—and direct government participation extending lines into rural
areas that were not commercially profitable. By the early twentieth
century, the network exceeded 40{,}000 kilometers, designed primarily
to connect the agricultural hinterland of the Pampas, the Northwest,
and the Northeast to the port of Buenos Aires. This radial,
export-oriented design was both cause and consequence of the
agricultural export boom of the first globalization: from 1869 to
1914, Argentine real exports grew approximately 500 percent
\citep{fajgelbaumredding2022}. By the start of the First World War,
Argentina was the eighth richest country in the world.

Starting in the 1930s, export growth slowed considerably and profit
margins of railroad companies declined. Investment in railroad
infrastructure stalled, and the aging network experienced rising
operational costs. Railroads also began to face competition from the
automotive industry—Ford produced the first Model~T in Latin America
in Buenos Aires in 1925. In 1932, Law 11.658 created a National Trunk
Road System and a lump-sum fuel tax whose proceeds were allocated
entirely to road construction and improvement; by 1935, there were
almost 3{,}000 kilometers of paved roads. During the Second World War,
Argentina remained neutral and its young industrial sector flourished
while factories in the developed world were retooled for war
production. Railroads experienced a ``second life'' as wartime food
shortages drove up international prices for Argentine agricultural
exports, helping the country accumulate reserves of \$1.6~billion
\citep{whitaker1954}. But by 1945, industry's contribution to GDP
surpassed that of agriculture for the first time \citep{potash1969}.

Supported by unions, the military, and the new industrial
bourgeoisie, General Juan Domingo Perón reached the presidency in
1946. His economic policy rested on three pillars: strengthening the
domestic market, centralized planning with state ownership of
strategic sectors, and industrial development
\citep{whitaker1954, pahowka2005}. The showcase of this program was
the acquisition and nationalization of the private railroad companies
in 1947 for \$600~million—a purchase widely debated at the time, with
economists including Raúl Prebisch arguing that the price was
excessive given the crumbling state of the infrastructure
\citep{rock1987}. By the early 1950s, when the fortuitous advantages
of the war years had faded, Argentina's economy slipped into
depression. Perón moved to protect nascent industry, armoring the
automotive and petrochemical sectors against international
competition. The road network continued to grow at a slow, steady
pace, reaching 9{,}000 kilometers of paved roads by 1955. The railroad
network, meanwhile, remained stable at around 43{,}500 kilometers, but
with a shrinking budget it suffered from poor maintenance and aging
equipment.

In 1955, a coup removed Perón from government, and the country
endured several years of political and economic instability. In 1958,
Arturo Frondizi was elected president in an election from which Perón
was banned. Frondizi's economic plan centered on fostering
industrialization through foreign direct investment: Decree 3693 of
1959 promoted investment in the national automotive industry, and 23
new factories were approved within a year. His administration was
equally aggressive in promoting road construction. In 1958, the
lump-sum fuel tax was replaced by a 35 percent ad valorem tax
(Decreto Nacional 505/1958), and in 1960 Law 15.274 instituted a
sales tax on vehicles—in both cases, funds were permanently allocated
to road construction and improvement. As for railroads, Frondizi
initially sought to modernize the crumbling infrastructure, but the
investment required to achieve a modern network while maintaining its
current size was nearly prohibitive \citep{lopez_waddell2007}. By
1960, annual railroad losses had reached approximately \$280~million
per year, accounting for almost 80~percent of the federal government
deficit \citep{keeling1993}.

\subsection{The Larkin Plan (1960--1962)}
\label{ssec:larkin}

In this context, the Argentine government agreed with the World Bank
to commission a team of experts to evaluate the long-term economic
and strategic viability of the national transport infrastructure. The
team was led by General Thomas Larkin, an American expert in military
logistics who had been decorated for his efforts supplying combat
troops during the Second World War and had experience consulting on
railroad modernization in France, Germany, and Japan. Between 1959
and 1961, the World Bank mission scrutinized the profitability of the
existing railroad and road networks. The resulting study, published
by the Public Works Ministry in 1961 as \emph{Transportes
Argentinos: Plan de Largo Alcance} (``Argentine Transportation:
Long-Run Plan''), comprised three volumes totaling approximately
3{,}000 pages. The first volume contained an executive summary with
diagnosis, methodology, results, and recommendations organized by
transport mode and time horizon; the remaining two volumes, of
approximately 1{,}000 pages each, provided detailed
analysis.\footnote{We accessed physical copies of the plan at the
Mariano Moreno National Library in Buenos Aires. We scanned and
digitized two key objects: a map describing the operational railroad
network in 1960 (our baseline) and three tables describing the
diagnosis and recommendations for each railroad segment.}

In a world of cheap and stable oil prices, and given the poor
material and financial condition of the railroad system, the plan's
main recommendations were dramatic: closure of approximately
15{,}000 kilometers of railroad track (roughly 32~percent of the
existing network), dismissal of around 70{,}000 railroad employees
(from a total of approximately 200{,}000), disposal of the entire
fleet of steam vehicles, closure of most national railroad workshops
and factories, and aggressive improvement and construction of paved
roads—in some cases to substitute for closed railroad segments, but
also to create new connections in the network \citep{larkin1962}. The
plan proposed the construction of approximately 10{,}000 kilometers of
paved roads over ten years.

A key feature of the Larkin Plan's design is that it did not study
the entire railroad network with equal intensity. Due to budget
constraints and the vast size of the existing network, only railroad
segments that fell below specific freight density thresholds were
``studied'' in detail—meaning that the experts checked each
segment's explicit revenues and costs and analyzed whether closure
was warranted. The thresholds were 1~million gross tons per year of
cargo passing through a given segment, or 500 tons per kilometer per
year of cargo entering the network through a given segment. Only
39.6~percent of the railroad network was studied. For each studied
segment, the experts issued one of three recommendations: maintain,
close, or perform a new study. Crucially, the likelihood of the plan
recommending closure was much higher for studied segments than for
nonstudied ones—a discontinuity that we exploit for identification.

The plan's recommendations were announced by President Frondizi in
1961, but massive protests and strikes erupted in many cities. After
several months of clashes and social unrest, the government decided
to abandon the plan. In the short term, only approximately
1{,}000 kilometers of track were closed. However, despite being
abandoned after just a few months, the Larkin Plan was the only
long-term comprehensive study of the national transport
infrastructure for more than 30 years. It is therefore likely that
its recommendations influenced national transport policy for decades
ahead. As we document below, railroad closures in the 1960s and
1970s were approximately three times more likely for segments that
had been studied by the plan, probably because no new large-scale
study was ever undertaken to supersede it.

\subsection{Implementation: Rail Closures and Road Expansion,
            1960--1991}
\label{ssec:implementation}

The Larkin Plan marked a trend break in the evolution of both
transport networks (Figure~\ref{fig:networks}). For railroads, it
initiated a long-term disinvestment and dismantling process. For
roads, it started an almost two-decade period of aggressive
construction and improvement. These large-scale changes occurred in
a context of population growth (from 20 to 33 million between 1960
and 1991) and increasing GDP per capita (growing approximately
35~percent at an average annual rate of 0.95~percent).

Railroad closures proceeded in two distinct waves. The first wave,
from 1960 to 1966, closed approximately 4{,}000 kilometers of track
before political opposition brought the process to a halt. The
second wave came under the military dictatorship that seized power
in 1976: without union resistance, the military government ordered
the immediate closure of over 6{,}000 additional kilometers of track
between 1976 and 1979. After 1979, the network stabilized—no further
closures or new construction occurred between 1979 and 1986. In
total, the railroad network fell from approximately
43{,}500 kilometers in 1960 to approximately 33{,}500 kilometers by
1986, a reduction of roughly 23~percent.\footnote{We end our sample
period in 1991 because there was a census in that year and because
we wish to abstract from the potentially different effects of the
privatizations of large fractions of the railroad and road networks
during the 1990s.}

Road expansion, by contrast, was continuous throughout the period.
Fueled by the dedicated fuel and vehicle taxes established under
Frondizi, the paved road network grew from approximately
10{,}000 kilometers in 1954 to approximately 28{,}000 kilometers by
1986—an increase of roughly 180~percent. Gravel roads also expanded,
though more modestly. By 1986, the combined paved and gravel road
network had reached approximately 35{,}000 kilometers, roughly equal
to the remaining railroad network. Importantly, the spatial pattern
of road expansion did not simply mirror the pattern of railroad
closures: roads were not systematically built to replace closed rail
lines. As we show in Section~\ref{sec:data}, the correlation between
rail closures and road expansion at the district level is weak,
meaning that the two shocks affected largely different sets of
districts.

\subsection{Scale Economies and Sector-Mode Specialization:
            Motivation and Hypotheses}
\label{ssec:scale_economies}

A feature of the Argentine setting that motivates part of our
analysis is that railroads and roads may serve different economic
sectors, owing to fundamental differences in their cost structures.
Railroads exhibit strong scale economies: high fixed costs (track
maintenance, stations, rolling stock) but low variable costs per
ton-kilometer at high traffic density. Roads, by contrast, have more
constant returns to scale—per-unit costs decline less steeply with
volume. The Larkin Plan itself was designed around this logic: it
studied segments precisely because their traffic density was too low
for rail's scale economies to be realized.

Data from \citet{baumgartnerpalazzo1969}, who estimated transport
costs for Argentina using detailed operational data, are consistent
with this pattern. At high cargo density (1{,}000 tons per day),
railroad costs were 0.465 pesos per ton-kilometer compared to 1.000
for roads—rail was more than twice as cheap. At low cargo density
(100 tons per day), the relationship reversed: railroad costs rose
to 13.490 pesos per ton-kilometer while road costs were 8.266—road
was substantially cheaper (Table~\ref{tab:cost_params}). If these
cost differences translated into actual shipping patterns, they
would imply a natural sector-mode specialization: agricultural
products—bulky, heavy, shipped in large volumes—would fall in the
high-density regime where rail has a cost advantage, while
manufactured goods—lighter, more varied, shipped in smaller
batches—would fall in the low-density regime where road is cheaper.

% [PLACEHOLDER: 1 paragraph on freight evidence.
%  (1) Data from Santa Fe Railroad showing manufactured goods faced
%      higher effective freights than agricultural goods.
%  (2) As road network expanded, railroads specialized in bulk
%      products — share of bulk cargo increased over time.
%  (3) Cite slides Figures 2–5 for evidence.
%  This paragraph depends on having access to the freight data
%  figures. Frame as suggestive, not conclusive.]

Whether these cost differences generate meaningful mode-specific
effects on regional economic outcomes is ultimately an empirical
question. It is not obvious a priori that they do: if road
expansion sufficiently compensated for rail closures, or if firms
and workers adjusted quickly to the new transport configuration,
the sectoral composition of economic activity might not have
changed. We treat the existence of mode-specific effects as a
hypothesis to be tested in Sections~\ref{sec:results} and~\ref{sec:counterfactuals},
rather than as an established fact. The cost data presented here
motivate the hypothesis but do not settle it.
